\chapter[World War II Becomes Global]{How the Second World War Became Truly Global}
\section{A Can of Meat That Traveled the World}
First, pick up an object. It weighs a little over three hundred grams: a rectangular tin. Lift the key-shaped metal tab on top and turn it several times. Inside lies tightly pressed pink cooked meat, edged with translucent jelly.

This is Spam, American canned luncheon meat introduced in 1937 as inexpensive protein for poor households during the Depression. Travel around the world in 1945, however, and you would encounter it in astonishing places.

In London, thin slices entered air-raid shelters, food aid shipped under American Lend-Lease. In Red Army trenches, soldiers pried open identical cans with bayonets, delivered by American freighters and trucks along Arctic routes or Iranian roads. On Guadalcanal, American soldiers carried it through rainforest and ate it until sick of it. In Hawaii, Okinawa, and Philippine villages, islanders traded with troops for cans and eventually incorporated them into everyday cooking.

Pause over the strangeness. London shelters, snowy Soviet trenches, and Pacific islands belonged, on contemporary maps, to three "different wars": the European war, the Soviet-German war, and the Pacific war. Military textbooks do teach them separately. Yet an unremarkable can joined them on one production line. American factories made it, American freighters carried it, and strangers on separate fronts opened identical tins in the same month.

Here is our question: how did the Second World War become a genuinely "world" war? It began as geographically separate conflicts, China against Japan in Asia, Germany against Britain and France in Europe. What welded them together, drawing South American rubber, Middle Eastern oil, Indian soldiers, Chinese peasants, and Pacific villages into one machine? How did that machine then change the world we still inhabit?

\section{A Summer Night in Chongqing, 1941}
Time: a night in late June 1941. Place: an air-raid shelter halfway up a hillside in Chongqing, Sichuan.

Chongqing was China's wartime capital. Japanese forces had occupied Wuhan but could not penetrate to this city among mountains and the Three Gorges, so turned to aircraft. From 1938, bombers repeatedly arrived, favoring clear days and moonlit nights. Every household kept an "alarm bag" of dry food, water, and essential documents, ready to run at the siren.

Tonight over a hundred people crowd the shelter. Cut into rock, it is deep, narrow, damp; water seeps down walls and drips onto necks. A few oil lamps hang on the stone, flames swaying in the breath of a hundred people. Sweat, kerosene, and sour diapers fill the air. A baby cries; its mother puts her nipple in its mouth, turning cries to muffled whimpers. Outside, distant thunderlike thuds arrive one after another: bombs across the city. Nobody talks. Everyone listens, judging the distance.

Deep inside, a middle-aged man in glasses suddenly speaks, softly but clearly. Neighbors call him Mr. Zhou. Before the war he taught Chinese in a secondary school; flight left only him and a nephew from his family.

"The Germans have invaded the Soviet Union. June twenty-second. Millions of men, thousands of tanks, pushing from Poland into Ukraine."

A moment's silence, then murmurs. A porter-looking man asks, "Which way is the Soviet Union? What's it to us?" Zhou answers, "North, next to us. If Germany defeats it, they will have their hands free..." Another interrupts: "Won't that please the Japanese? Didn't Germany sign an alliance with them?"

Someone else gives a bitter laugh. "What have European wars got to do with us? We've been fighting ten years. Who cared about us?"

The words land heavily; nobody answers. Ten years since the explosion outside Shenyang in 1931 and the Northeast's occupation; four years of full-scale war since the Marco Polo Bridge incident in 1937. Everyone remembers three weeks earlier. On the evening of June 5, an exceptionally long air raid drove far more people than capacity allowed into Jiaochangkou's great shelter tunnel. Doors closed, ventilation failed, and hundreds or thousands suffocated or were crushed during hours of bombing. Estimates still range from hundreds to thousands without agreement.

They do not know that summer 1941 is a hinge of world history. The war dropping bombs above them and the European war on Zhou's radio are rapidly converging. Within six months they will be welded at a place thousands of kilometers away called Pearl Harbor. Aircraft bombing Chongqing burn imported high-octane fuel. A few months later, one oil embargo will bind those aircraft's fate to the American Pacific Fleet's.

\section{Which Year Did the War Begin?}
Before going on, consider a question worth disputing seriously: in which year did the Second World War break out?

Your textbook probably says September 1, 1939: Germany invaded Poland, Britain and France declared war, and the Second World War began. This answer spread widest because the earliest writers of "world history" largely sat in European and American studies.

Ask someone in that Chongqing shelter. China's war began on September 18, 1931, when Japan's Kwantung Army blew up a small stretch of the South Manchuria Railway near Shenyang, blamed Chinese troops, occupied the city, and swallowed the entire Northeast within months. Or they might date full-scale war to the Marco Polo Bridge incident of July 7, 1937. After that, no Chinese family remained wholly outside it.

Ask a Pole: 1939. An Ethiopian: 1935, when Italian forces invaded the ancient Horn of Africa kingdom, bombed undefended towns, and even used poison gas. The Western-dominated League of Nations offered little beyond condemnation and halfhearted sanctions. Many later thought this success taught Hitler that aggression carried no cost. An American: December 7, 1941, Pearl Harbor. A Soviet citizen: June 22, 1941.

Five honest witnesses offer five starting dates and cannot persuade one another. Nobody has simply forgotten. The deeper question is \textbf{who draws the boundaries of "a war"?} Starting on September 1, 1939, implicitly makes European gunfire the real war and northeastern China's occupation, Nanjing's massacre, and Ethiopian poison gas merely a "prelude." Starting in 1931 feels distant to European readers.

More difficult than the starting date is the welding. Between 1937 and 1941, the Asian and European conflicts really were separate: Japan did not help Germany attack France, nor Germany help Japan attack Chongqing. Soldiers did not meet and commands were unrelated. Two hammer blows in the latter part of 1941 joined them. In June, Germany broke its treaty and invaded the Soviet Union, expanding European war eastward across the continent. In December, Japan attacked Pearl Harbor and simultaneously launched broad offensives across Southeast Asia and the Pacific. America declared war; Germany and Italy then declared war on America. Around Christmas 1941, all major powers were involved, with fighting across Chinese plains, Soviet snowfields, North African deserts, Atlantic shipping routes, and Pacific islands.

That war grew is easy to understand. Harder is \textbf{why it had to grow}. Why could no state remain on its own battlefield and keep war "local"? The shelter's questioner was wrong that Europe's fighting had nothing to do with him, as he would discover within months. To explain how wrong, we need more voices and a new tool.

\section[Soldiers under Imperial Flags]{Fighting for Whom? Soldiers under Imperial Flags}
Change perspective: look not at great powers on maps but at people pressed beneath the maps.

Begin with India. British India's army grew to about 2.5 million during the war, one of history's largest volunteer armies. Indian soldiers fought Rommel's Afrika Korps in North African deserts, contested Italian mountain villages, and battled Japanese troops in Burma's jungles. India was not a belligerent state; nobody asked Indians whether to fight. The British colonial government declared war for them.

Try a serious exercise: \textbf{state both sides' reasons so fully that each side would recognize its own position.}

Speak for Indian recruits. First, pay was real. For a poor rural family, a soldier son meant food, and death benefits were written commitments. Second, Japan reached India's doorstep: in 1944 Japanese troops invaded the Imphal area of the northeast. Home defense was no abstraction. Third, many sincerely saw a war against racism and aggression and believed Britain would honor promises of Indian freedom afterward. None of these reasons was pretended.

Now the other side. Indian politician Subhas Chandra Bose chose the opposite route. Escaping British surveillance, he reached Japanese-occupied Singapore and reorganized captured Indian soldiers into the Indian National Army. His slogan broadly promised, "Give me blood and I will give you freedom." Alongside Japan, it fought toward India's frontier. British voices called him traitor and Japanese puppet. Give his reasoning fully: Congress sought a British commitment to postwar independence before cooperation, and Britain refused. If Britain fought for "freedom" while denying India's, "Britain's difficulty is India's opportunity." Why not use Japanese help to expel the British? Most soldiers in this army were not Fascist believers; they wanted colonizers gone.

You need not choose a side. See instead that, for colonized peoples, the war was never simply "justice or evil" but "whose chains replace whose?" Two Indian armies fired at one another in Burma, each believing its men died for India's freedom.

Now Africa. French colonial forces included the famous Senegalese tirailleurs, a general name for men from across French West Africa. They bled for metropolitan France in 1940. After defeat, German troops shot groups of captured Black soldiers, deeming them unworthy of prisoner-of-war status. When the Allies returned to France in 1944, Free French forces contained many Africans, but before Paris's liberation French commanders replaced them at the front with Whites, preferring "White French faces" in the liberation picture. The sharpest blow came after the war: in late 1944, at Thiaroye camp near Dakar, Senegal, French troops fired on returning African former prisoners demanding overdue pay. Dozens fell; the precise toll remains disputed. Four years fighting for "Free France" brought bullets at their home camp's gate.

Consider the Atlantic Charter too. In August 1941, Roosevelt and Churchill signed this document aboard warships off Newfoundland. Its third clause declared peoples' right to choose their government. Nationalists across India, Vietnam, Indonesia, and African colonies greeted it as a check payable to them. Months later, Churchill explicitly told Parliament the principle applied only to Nazi-occupied European countries, not the British Empire. The check acquired a restricted-use clause.

An easily misunderstood counterexample: Asian peoples did not all recognize Japanese aggression immediately. Japan expanded under the banner of a "Greater East Asia Co-Prosperity Sphere," liberating Asia from White colonial rule. Some nationalists in Burma, Indonesia, and the Philippines initially welcomed a force expelling British, Dutch, and American rulers. Aung San, later Burma's founding father and Aung San Suu Kyi's father, helped organize an army with Japan. Why did propaganda persuade? Colonial grievances were real; "Asia for Asians" touched an actual wound. Welcomers soon found new rulers harsher than old: forced labor, seized food, casual massacre. Hundreds of thousands of conscripted Javanese laborers died on projects including the Burma Railway. Aung San later turned his army against Japan. The lesson: \textbf{a slogan is often dangerous not because it is entirely false but because it is seventy percent true}. Judge a "liberator" by labor camps, not banners.

Finally, the Pacific. Coastwatchers in the Solomon Islands, many local islanders, risked summary execution to radio Japanese movements and saved countless Allied lives. Papuan New Guinean carriers hauled wounded men and ammunition along muddy Owen Stanley Range trails; grateful Australians called them "Fuzzy Wuzzy Angels." Their names rarely entered national commemorative books afterward. China had millions without uniforms beyond its regular and guerrilla fronts: two hundred thousand Yunnan laborers who hammered and carried their way through more than a thousand kilometers of mountain road in nine months to build the Burma Road; workers forcibly sent to northeastern and northern China or Japanese mines; and tens of thousands of women from China, Korea, and Southeast Asia forced into the "comfort women" system. Their experiences went largely unrecorded for decades until survivors testified.

Together these voices reveal something a Europe-centered story cannot contain: much of the war's manpower, food, raw materials, and dead came from people deemed unqualified to participate as nations.

\section[Shipping Routes and Ledgers]{A Bridge for Thought: Shipping Routes and Ledgers}
War books favor campaign maps, arrows from Warsaw to Paris or Moscow to Berlin. Useful though they are, they can hide why war became global. Try a "four-question source-tracing method": \textbf{for any wartime event, set aside front lines and ask where the oil, food, soldiers, and money and weapons come from}. Draw lines along the answers, and you produce a network rather than a map.

Begin with the hardest question: why attack Pearl Harbor? Striking America's industrial giant looked nearly mad militarily; Yamamoto Isoroku himself knew a long war was unwinnable. Trace resources and the chain appears. First, oil: Japan was oil-poor, importing about eighty percent before the war, mostly from America. In July 1941, protesting occupation of southern French Indochina, America joined Britain and the Netherlands in freezing Japanese assets and imposing an oil embargo. Stocks could sustain the war machine for little more than a year. Second, alternatives: Dutch East Indies oilfields, now Indonesia, and Malayan rubber and tin. Resources for invading China lay in Southeast Asia. Third, obstacles: British, American, and Dutch colonies meant southern expansion required war with Britain and America. American troops held the Philippines, the Pacific Fleet Pearl Harbor. Unless disabled, they exposed the invasion convoys' flank. Thus: war in China needs oil$\rightarrow$America cuts oil$\rightarrow$seizing oil requires southern expansion$\rightarrow$expansion requires striking America first$\rightarrow$Pearl Harbor. An Asian war necessarily flowed through oil pipes into the Pacific.

Now Germany. Ideology publicly justified invading the Soviet Union, but resource accounting was equally clear: Ukraine's black soil was a granary, and Baku's Caucasian oilfields were what Germany's war machine most craved. The principal direction of its 1942 summer offensive was the Caucasus. Germany depended heavily on Romania's Ploiești oilfields, drawing Allied bombing routes along the oil pipes into Romanian skies.

The Allied network was more spectacular still. Lend-Lease made American factories the "arsenal of democracy." Tanks, trucks, cans, and boots sailed Arctic routes to Murmansk, around the Cape of Good Hope to India, and through the Red Sea to Egypt. After Japan cut the Burma Road, Chinese and American pilots opened the "Hump" air route across Himalayan ridges from India to Kunming. Hundreds of aircraft and more than a thousand crew members were lost over several years. Aluminum wreckage glittering in snowy valleys earned the name "Aluminum Valley." In the Atlantic, German submarines hunted convoys while Allied escorts, long-range aircraft, and codebreaking fought them throughout the war. Everyone understood: cut the shipping routes and lose the war.

The tool explains why \textbf{no country could remain outside}. In the industrial world of the 1940s, oil, rubber, food, and shipping already stitched major economies into a network. War at a major node necessarily spread along its connections as each side attacked enemy supplies and protected its own. Globalization was not war's result but its transmission mechanism. Try these questions with any supposedly "local" war in the news; the network will probably already cross national borders.

\section{Promises on Two Sheets of Paper}
Now read primary material. Midway through the war, future victors were already arranging the postwar world on paper. Those words reveal more than battlefield bulletins about why the war was fought and what world it produced.

First, the Cairo Declaration of November 1943, issued after Chinese, American, and British leaders met. Its passage on China stated:

\begin{quote}
"The purpose of the three countries... is that territories Japan stole from China, including the four northeastern provinces, Taiwan, and the Penghu Islands, be returned to the Republic of China."
\end{quote}
(Cairo Declaration, issued December 1, 1943.)

Weigh this document. Since 1931, China had borne Japanese aggression alone. At Cairo it first sat at the table as one of the Four Powers, with restoration of lost territories written into their common statement. For a country fighting over a decade with tens of millions dead, these lines were bought in blood. Notice the phrasing too: it addressed "Chinese territories stolen by Japan," an anticolonial list stopping at what Japan had taken before 1914, without mentioning British or American Asian colonies.

Second, the United Nations Charter of 1945. Before war fully ended, fifty countries' representatives drafted a new international organization's constitution in San Francisco. Its preamble opened:

\begin{quote}
"We the peoples of the United Nations are determined to spare future generations the scourge of war, which twice in our lifetime has brought humanity suffering beyond description."
\end{quote}
(Preamble to the United Nations Charter, signed June 26, 1945.)

"Twice": only twenty-one years after the First World War ended, humanity began a greater one. The sentence contains real fear and ambition: install a permanent brake on the world.

Read these documents with the Atlantic Charter and a pattern appears. Wartime promises are a special currency, required to do contradictory things: promise citizens and allies a new world worth dying for---self-determination, freedom from war, restored territory---while preserving signatories' old world of empires, influence, and military bases. When Roosevelt and Churchill wrote that peoples could choose governments, one imagined undermining old colonial systems and opening space for American influence; the other imagined limiting it to Europe. Chinese representatives signed the UN Charter thinking "never again," while Indian, Vietnamese, and Indonesian representatives were absent, still colonized.

How much was redeemed? China recovered Taiwan and the Northeast and became a permanent Security Council member; the UN still functions. The other side of the check: British, Dutch, and French forces returned to Malaya, Indonesia, and Indochina to reconnect colonial rule. But it would not reconnect. War had armed and mobilized Asian and African peoples and educated them in "national self-determination." Over the next two or three decades, independence wars redeemed that discounted promise with interest. One of the war's deepest consequences was to hollow out every empire's foundations while crushing Germany's, Japan's, and Italy's.

\section[Thorough Killing: Three Explanations]{Why Was the Killing So Thorough? Three Explanations}
Approach the chapter's heaviest subject by establishing facts before explanations.

The death toll was unprecedented: estimates range from sixty to eighty-five million. Uncertainty itself reflects the catastrophe. China and the Soviet Union each lost tens of millions; entire villages disappeared unregistered. Civilian deaths exceeded military deaths. In the six weeks after Nanjing fell, Japanese troops committed mass murder, rape, and looting, killing hundreds of thousands. China's recognized figure is three hundred thousand; the International Military Tribunal for the Far East found more than two hundred thousand. Nazi Germany conducted planned, industrialized, continent-wide extermination, murdering roughly six million Jews and systematically killing hundreds of thousands of Roma, formerly called Gypsies, disabled people, and others. More than three million Soviet prisoners of war died from hunger, cold, and execution. Late in the war, over seven million foreign civilian laborers worked under compulsion in Germany alone. Allied bombs also struck cities. Tokyo's overnight raid beginning March 9, 1945, burned about one hundred thousand people to death, mostly civilians; Dresden's February firestorm killed about twenty-five thousand. Atomic bombs destroyed Hiroshima and Nagasaki on August 6 and 9, with combined deaths reaching roughly two hundred thousand by year's end.

Why? Separate four things: what happened, the facts above; why, our coming debate; how contemporaries understood it; and how we judge it. At least three serious explanations address the thoroughness of killing.

\textbf{First: ideology kills.} This account centers Nazi racism, Japanese militarism's "holy war," and extreme nationalism. Once an ideology defines people as "subhuman," "vermin," or an "inferior race," killing becomes sanitation. It explains extermination's most horrifying qualities, selectivity and systematic organization. Gas chambers, assembly-line killing centers, and bureaucratic victim lists were not merely "war out of control" but cool target management. One detail is especially telling: in 1944, with Germany losing in the east and transport desperately scarce, the regime still diverted valuable railway cars to carry Hungarian Jews to Auschwitz. Extermination outranked military need. Only ideology explains that. But it does not explain the Allies: Tokyo and Dresden's incendiaries and Hiroshima's atomic bomb also killed civilians en masse, although their users were not extermination regimes.

\textbf{Second: total war's industrial logic.} Shift from "who hates whom" to "what modern war depends on." Twentieth-century war depended on factories, railways, and fields. Behind one shell stood mines, steelworks, lathes, and farmers. Civilians became central components of war economies rather than bystanders. On this logic, destroying urban workers achieved what destroying an arsenal did. Bomber technology could not deliver precision, making "night area bombing" standard. This explains scale and indiscrimination, why democracies planned the destruction of cities and warfare became increasingly "everyone's." Two limits remain. It risks technological determinism: repeated claims that bombers and structures "forced" action quietly erase decision-makers. And it does not explain the Holocaust: exterminating Jews did not serve Germany's war economy and even harmed it. Structure explains carpet bombing, not gas chambers.

\textbf{Third: imperial violence returns home.} Thinkers including Martinique's Aimé Césaire, later author of Discourse on Colonialism, identify an uncomfortable genealogy. Concentration camps, racial laws, forced labor, and mass killing of "inferior peoples" had been rehearsed for decades in colonies before appearing in Europe. Hitler, in one sense, applied to Europeans what Europeans had long inflicted on people of color. This explains why "civilized world" narratives long treated Asian and African losses lightly. Nanjing's dead, Javanese forced laborers, and the millions who starved in Bengal in 1943, with wartime grain policies bearing responsibility, counted as lower-ranked casualties in a racial hierarchy. Its limit: genealogy is not identity. The Holocaust's ambition to erase an entire people has a distinctive, irreducible thoroughness, not simply "continued colonialism."

Each explanation grasps some truth but not all. Atomic bombing is their fiercest battleground. The standard account says it forced surrender and avoided potentially million-scale casualties invading Japan. That calculation was real: Japanese preparations for homeland battle were substantial, and Okinawa's horror frightened American planners. Another account also has evidence: American intelligence knew Japan was exploring Soviet mediation, and atomic bombs simultaneously signaled Stalin. The next confrontation had begun before this war ended. The Manhattan Project had also spent more than two billion dollars and produced bombs; established institutions had inertia. Historians still debate it. Take this methodological lesson: \textbf{when a major decision is allowed only one explanation, be wary}.

Before leaving, be clear: comparing explanations does not blur everything into "everyone was equally bad." Nazi extermination camps and Allied strategic bombing differed fundamentally in motive, targets, scale, and organization. Identifying Allied killing of civilians does not reduce Fascist guilt, just as recognizing British and French imperialism does not justify Japanese aggression. This distinction is the entrance requirement for discussing the war.

\section[How Society Makes Memory]{A Deeper Challenge: How Society Makes Memory}
After fighting ended, a quieter battlefield opened. The same war looks almost like different wars in different countries' monuments, textbooks, and films.

Introduce a sociological tool. French sociologist Maurice Halbwachs developed "collective memory" in the first half of the twentieth century; he himself died at Buchenwald in 1945. His central insight was broadly that memory is a building site, not a warehouse. Every recollection reconstructs the past according to present social needs. What a group remembers, forgets, or puts at the center of a public square depends on what it currently needs to become.

Scan national war memories through this theory and contrasting construction sites appear.

\textbf{America}'s main story is "the most just war": Pearl Harbor, Normandy's beaches, liberated camps, atomic bombs that "ended the war." It supports America's postwar "world leader" identity while dimming collective incarceration of Japanese American citizens and civilian costs of bombing German and Japanese cities.

\textbf{Russia}'s main story is the "Great Patriotic War": twenty-seven million dead, Stalingrad, Berlin's capture, and May 9 Victory Day, still its grandest national celebration. It remembers unprecedented sacrifice and victory, long leaving unspoken the prewar Soviet-German nonaggression pact and mass graves of Polish officers at Katyn.

\textbf{China}'s main story is fourteen years of resistance, from September 18 to victory, with the Nanjing Massacre central to national trauma. Since 2014, December 13 has been the National Memorial Day for its victims. This version changes too: replacing "eight-year resistance" with "fourteen-year resistance" was public reconstruction of memory.

\textbf{Japan}'s memory is most divided. One line emphasizes victimhood: burned Tokyo, Hiroshima and Nagasaki's mushroom clouds. Japan alone has suffered atomic bombing; opposition to nuclear weapons approaches national consensus. Another emphasizes perpetration: Nanjing, "comfort women," forced labor. This remains contested domestically, from prime ministers' Yasukuni visits to textbook descriptions of invasion, each affecting East Asia. Two memories struggle inside one country while neighbors see only one wall of its building site.

\textbf{Germany} is the most distinctive case, undergoing decades of open memory reconstruction. Early postwar Germans favored victim stories: bombed cities, flight from the east, hunger, with widespread silence on the Holocaust. Generations deliberately changed this. In 1970, Chancellor Willy Brandt knelt at Warsaw's ghetto memorial. In 1985, President Richard von Weizsäcker called May 8, Germany's surrender date, a "day of liberation": liberation from tyranny, not defeat. Notice the operation. It did not deny defeat but rearranged its meaning, allowing a nation to face its past differently. This is a textbook scene for Halbwachs's theory.

Another persistent blind spot is that \textbf{colonial memories enter none of the principal narratives}. The 2.5 million Indian soldiers, Thiaroye veterans, Pacific carriers, and Bengal famine victims were long absent from European and American commemorative systems, only gradually recovered by scholars and descendants in recent decades. Whose sacrifice is carved in stone, whose confined to a statistical table? This is memory politics' most revealing test.

The theory offers a new way to read news. When a country changes textbooks, creates a memorial day, or erects or removes monuments, first ask whether it fits historical facts; that question is essential. Then ask what kind of country it now wants to become. Memory disputes are always disputes about present identity.

\section{Echoes Still Sound Today}
This war of more than eighty years ago still operates in your world, in altered forms.

It left the UN's skeleton. The Security Council's permanent five---China, America, Russia as Soviet successor, Britain, France---are the wartime victors' list, freezing 1945's power arrangement into today. You can see achievements in peacekeeping, mediation, and humanitarian relief, and paralysis when great-power vetoes block action. Decades of reform calls encounter the same difficulty: those needing reform hold the veto. Here is a first unresolved inheritance: victors write the order, and victors grow old.

It left a human-rights vocabulary. The Universal Declaration of Human Rights and Genocide Convention of 1948 and new Geneva Conventions of 1949 grew directly from the war's dead. "Crimes against humanity" entered law through Nuremberg and eventually reached today's International Criminal Court. The revolutionary claim was that leaders and soldiers could not defend slaughter through "orders" or "sovereignty": humanity stood above sovereignty. A vast gap still separates ideals and enforcement, but without the vocabulary we could not even measure it.

It left the nuclear shadow. After Hiroshima, humanity invented the strange peace called "deterrence": nobody dares act first because action means mutual destruction. The Nuclear Non-Proliferation Treaty tries to lock the nuclear club, but the keyhole keeps widening. Your generation will grow up among more than ten thousand nuclear warheads.

It also left memory wars. Japanese textbook wording and officials' Yasukuni visits repeatedly provoke Chinese and Korean protests; Europeans likewise disagree over commemoration. This is not simply "digging up old accounts." Halbwachs's building site is still active, each side defining itself through memory.

Finally, the war may be in your schoolbag. Open history textbooks with peers in Japan, Korea, Russia, and America and you find five wars with similar outlines but radically different centers. When you encounter another version, do not immediately declare its reader "brainwashed." Ask what building their memory site is constructing, and what yours is constructing too.

\section[Remembrance: Dam or Reservoir?]{Your Question: Does Remembrance Build a Dam or Fill a Reservoir?}
Return to the Spam tin and the Chongqing shelter.

You have seen separate regional wars welded by oil, shipping, empire, and ideology into a global machine. After it passed, humanity swore "never again" through the UN and human-rights declarations and kept the dead through commemorations.

Here is the final question:

\textbf{Can remembering war prevent war?}

One side has strong reasons: a society forgetting genocide has no immunity; decades of German remembrance helped European reconciliation; without memorial days and museums, victims die a second death. The other is no weaker. Memory sites are never pure; they can nourish hatred and mobilize the next war. Serbian-Croatian and Israeli-Palestinian conflicts show how carefully preserved old wounds fuel conflict. Does a people teaching children only to "remember hatred" build a dam or store water for the next flood?

If remembrance is necessary, what is the right dose? Up to what point is it a vaccine, and beyond what point a virus? How should a country---China, Japan, or Germany today---remember the war eighty years ago while doing justice to the dead without wronging the living?

Give your answer and reasons. Ideally, those reasons should hold together the can of meat, the shelter, and the knee bent before Warsaw's memorial.
